\section*{AXIOMS AND SCHEMES OF THE THEORY OF SETS}
\phantomsection
\addcontentsline{toc}{section}{Axioms and Schemes of the Theory of Sets}

S1. If \(A\) is a relation, (A or A) \(\Longrightarrow A\) is an axiom.

S2. If A and B are relations, the relation \(A \Longrightarrow (A \text{ or } B)\) is an axiom.

S3. If A and B are relations, the relation \((A \text{ or } B) \Longrightarrow (B \text{ or } A)\) is an axiom.

S4. If A, B, and C are relations, the relation

\[
(A \Rightarrow B) \Rightarrow ((C \text {   or   } A) \Rightarrow (C \text {   or   } B))
\]

is an axiom.

S5. If R is a relation, T a term, and x a letter, the relation \((T|x)R \Longrightarrow (\exists x)R\) is an axiom.

S6. Let \(\pmb{x}\) be a letter, \(\pmb{T}\) and \(\pmb{U}\) terms, and \(R\{\pmb{x}\}\) a relation. Then the relation

\[
(T = U) \Longrightarrow (R \{T \} \Longleftrightarrow R \{U \})
\]

is an axiom.

S7. Let R and S be relations and x a letter. Then the relation

\[
((\forall x) (R \leftrightarrow S)) \Longrightarrow (\tau_ {x} (R) = \tau_ {x} (S))
\]

is an axiom.

S8. Let \(R\) be a relation, \(x\) and \(y\) distinct letters, \(X\) and \(Y\) letters distinct from \(x\) and \(y\) which do not appear in \(R\). Then the relation

\[
(\forall y)(\exists X)(\forall x)(R \Longrightarrow (x \in X)) \Longrightarrow (\forall Y)\operatorname{Coll}_x((\exists y)((y \in Y) \text{ and } R))
\]

is an axiom.

A1. \((\forall x)(\forall y)(((x\subset y)\text{ and }(y\subset x))\Longrightarrow (x = y)).\)

A2. \((\forall x)(\forall y)\operatorname{Coll}_z((z = x)\text{ or }(z = y)).\)

A3. \((\forall x)(\forall x')(\forall y)(\forall y')(((x, y) = (x', y')) \Longrightarrow ((x = x')\text{ and }(y = y'))).\)

A4. \((\forall X)\operatorname{Coll}_{Y}(Y \subset X).\)

A5. There exists an infinite set.
